% chapters/02-literature-review.tex — บทที่ 2 การทบทวนวรรณกรรม
\chapter{การทบทวนวรรณกรรม}
\label{ch:lit}

\section{กระบวนการคำนวณต้นทุนงานก่อสร้าง}

การคำนวณต้นทุนงานก่อสร้างเป็นกระบวนการประมาณการค่าใช้จ่ายที่จำเป็นในการ
ดำเนินโครงการก่อสร้างให้แล้วเสร็จตามแบบและข้อกำหนด โดยอาศัยข้อมูลแบบก่อสร้าง
รายการปริมาณงาน (BOQ) ราคาวัสดุ ค่าแรง และค่าใช้จ่ายที่เกี่ยวข้อง.

\subsection{ต้นทุนวัสดุ}

ต้นทุนวัสดุ (Material Cost) คือค่าใช้จ่ายที่เกิดขึ้นจากการจัดหาวัสดุที่ใช้
เป็นองค์ประกอบหลักในการก่อสร้าง วัสดุก่อสร้างที่เกี่ยวข้องอาจประกอบด้วย
วัสดุโครงสร้าง เช่น ปูนซีเมนต์ เหล็กเสริม ทราย หิน รวมถึงวัสดุตกแต่ง
เช่น กระเบื้อง สี และอุปกรณ์ประกอบอาคารต่าง ๆ.

\section{หลักการคำนวณต้นทุนวัสดุต่อหน่วย}

\subsection{ต้นทุนต่อหน่วย}

ต้นทุนต่อหน่วย (Unit Cost) เป็นแนวคิดพื้นฐานทางด้านเศรษฐศาสตร์และการจัดการ
ต้นทุน ซึ่งถูกนำมาใช้อย่างแพร่หลายในการประมาณราคา การวางแผนงบประมาณ และ
การควบคุมต้นทุน \autocite{ashworth2015}.
ต้นทุนต่อหน่วยช่วยให้ผู้ประมาณราคาสามารถเปรียบเทียบต้นทุนระหว่างรายการงาน
หรือโครงการที่มีลักษณะใกล้เคียงกันได้อย่างเป็นระบบ.

\subsection{ขั้นตอนการคำนวณต้นทุนต่อหน่วย}

\begin{enumerate}
  \item \textbf{การสืบราคาวัสดุ (Material Price Investigation)} ---
        รวบรวมข้อมูลราคาวัสดุก่อสร้างจากแหล่งข้อมูลต่าง ๆ เช่น
        ราคากลางหน่วยงานภาครัฐ ผู้ผลิต ร้านค้าวัสดุก่อสร้าง หรือฐานข้อมูล
        ออนไลน์ (Modern Trade)
  \item \textbf{การคำนวณต้นทุนวัสดุต่อหน่วย} --- กำหนดราคาวัสดุต่อหน่วย
        พื้นฐาน เช่น ราคาต่อถุง ต่อกิโลกรัม หรือต่อลูกบาศก์เมตร
        จากนั้นนำมาพิจารณาร่วมกับปริมาณการใช้วัสดุตามแบบก่อสร้าง
  \item \textbf{การวิเคราะห์} --- เปรียบเทียบเชิงพื้นที่และเวลา รวมถึง
        ประเมินความสอดคล้องของราคา
\end{enumerate}

\subsection{ตัวอย่างการคำนวณต้นทุนต่อหน่วย}

\textbf{งานบุกระเบื้องผนังดินเผาเคลือบเซรามิคสีธรรมดา ขนาด 4"x4"
(ไม่รวมงานฉาบปูนรองพื้น)}

\begin{enumerate}
  \item กระเบื้องเคลือบเซรามิคสีธรรมดา 4"x4" จำนวน 1.10 ตร.ม.
        @ 419 บาท/ตร.ม. = $1.10 \times 419$ บาท/ตร.ม.
  \item ปูนกาวซีเมนต์ หนา 5 มม. จำนวน 5.25 กก.\ @ 3.14 บาท/กก.
        = $5.25 \times 3.14$ บาท/ตร.ม.
  \item ปูนยาแนว จำนวน 0.38 กก.\ @ 3.14 บาท/กก.
        = $0.38 \times 3.14$ บาท/ตร.ม.
  \item น้ำผสมปูน จำนวน 2.00 ลิตร @ 0.0164 บาท/ลิตร
        = $2 \times 0.0164$ บาท/ตร.ม.
  \item คิ้ว PVC (เล็ก) จำนวน 0.33 ม.\ @ 30 บาท/ม.
        = $0.33 \times 30$ บาท/ตร.ม.
\end{enumerate}

รวมงานบุกระเบื้องผนัง พื้นที่ 1 ตร.ม. = 522.33 บาท/ตร.ม.\\
ค่างานต้นทุน = ค่าวัสดุ + ค่าแรง = $522.33 + 201 = 723.33$ บาท/ตร.ม.

\section{ข้อจำกัดในการสืบค้นราคาวัสดุในปัจจุบัน}

\begin{enumerate}
  \item \textbf{ความกระจัดกระจายของแหล่งข้อมูลราคา} --- ราคาวัสดุเผยแพร่
        อยู่ในหลายแหล่ง ทำให้ผู้ประมาณราคาต้องสืบค้นจากหลายที่
  \item \textbf{รูปแบบข้อมูลไม่เป็นมาตรฐาน} --- หน่วยจำหน่าย ชื่อวัสดุ
        ขนาดสินค้า แตกต่างกันในแต่ละแหล่ง
  \item \textbf{ความแตกต่างของหน่วยจำหน่ายและหน่วยใช้งาน} --- เช่น
        ราคาเหล็ก 25{,}000 บาท/ตัน เทียบกับ 25 บาท/กก.
  \item \textbf{ความไม่ทันสมัยของข้อมูลราคา} --- ราคาทางการอาจเผยแพร่
        รายเดือนหรือรายไตรมาส
  \item \textbf{ความแตกต่างของราคาตามพื้นที่} --- ราคาวัสดุแตกต่างตามภูมิภาค
        ระยะทางขนส่ง และโครงสร้างตลาดท้องถิ่น
  \item \textbf{ข้อจำกัดในการเข้าถึงข้อมูลดิจิทัล} --- ข้อมูลบางส่วนอยู่ใน
        รูปแบบ PDF หรือรูปภาพ ทำให้ดึงข้อมูลเชิงโครงสร้างยาก
  \item \textbf{ความไม่สอดคล้องของชื่อและคุณลักษณะวัสดุ} --- วัสดุชนิดเดียวกัน
        อาจมีชื่อทางการค้าและขนาดต่างกัน
  \item \textbf{กระบวนการยังเป็นแบบกึ่งอัตโนมัติ} --- ผู้ประมาณราคายังต้อง
        สืบค้นเอง ใช้เวลามาก
\end{enumerate}

\section{เทคโนโลยีที่ใช้ในการดึงข้อมูล}

เทคโนโลยีที่ใช้ในการดึงข้อมูล (Data Extraction Technology) สามารถจำแนกได้เป็น
4 ประเภท ดังนี้

\subsection{Application Programming Interface (API)}

API เป็นเทคโนโลยีที่ใช้สำหรับการเชื่อมต่อและแลกเปลี่ยนข้อมูลระหว่างระบบ
คอมพิวเตอร์หรือแอปพลิเคชันต่าง ๆ ผ่านชุดคำสั่งหรือโปรโตคอลที่กำหนดไว้ล่วงหน้า
มักทำงานผ่านเครือข่ายอินเทอร์เน็ต โดยใช้รูปแบบการสื่อสารเช่น HTTP/HTTPS และ
ส่งข้อมูลในรูปแบบมาตรฐาน เช่น JSON หรือ XML.

\subsection{การสกัดข้อมูลด้วยปัญญาประดิษฐ์ (AI-Based Data Extraction)}

นำหลักการของปัญญาประดิษฐ์มาประยุกต์ใช้ในการดึงข้อมูลจากแหล่งข้อมูลที่หลากหลาย
ทั้งข้อมูลที่มีโครงสร้าง (Structured Data) และไม่มีโครงสร้าง (Unstructured)
โดยอาศัยเทคนิค Machine Learning, NLP, และ Pattern Recognition.

\subsection{Web Scraping}

Web Scraping เป็นกระบวนการสกัดข้อมูลที่แสดงผลบนหน้าเว็บไซต์ เช่น ข้อความ
ตาราง หรือราคา เพื่อนำมาจัดเก็บในรูปแบบที่สามารถนำไปประมวลผลได้อย่างเป็นระบบ
\autocite{rao2015}.
หลักการสำคัญคือการระบุแหล่งข้อมูล กำหนดรูปแบบเป้าหมาย และจัดเก็บใน
โครงสร้างที่เหมาะสม.

\subsection{Database Integration (DI)}

การบูรณาการฐานข้อมูล เป็นเทคโนโลยีที่ใช้ในการเชื่อมโยงและรวมข้อมูลจาก
ฐานข้อมูลหลายแหล่งให้สามารถทำงานร่วมกันเป็นระบบเดียว ผ่านเครื่องมือต่าง ๆ
เช่น Middleware, Data Warehouse, หรือ ETL.

\section{งานวิจัยที่ผ่านมาในอดีต}

\textcite{rao2015} แสดงให้เห็นถึงศักยภาพของการใช้เทคนิค Web Scraping ในการ
รวบรวมข้อมูลราคาจากแหล่งข้อมูลออนไลน์หลายแหล่ง ซึ่งช่วยแก้ไขปัญหาความล่าช้า
และความไม่ต่อเนื่องของข้อมูลราคา อย่างไรก็ตาม งานวิจัยดังกล่าวยังไม่ได้
เชื่อมโยงข้อมูลราคาเข้าสู่กระบวนการคำนวณต้นทุนต่อหน่วยในบริบทของงานก่อสร้าง
โดยตรง.

\textcite{chunyaem2019} ได้นำเสนอแนวทางการประยุกต์ใช้ปัญญาประดิษฐ์เพื่อ
จัดจำแนกและจัดโครงสร้างข้อมูลวัสดุก่อสร้างโดยอัตโนมัติ ซึ่งช่วยลดปัญหา
ความไม่เป็นมาตรฐานของข้อมูลนำเข้า แต่งานวิจัยดังกล่าวมุ่งเน้นไปที่ขั้นตอน
การจัดประเภทข้อมูลเป็นหลัก.

\textcite{liu2024} ได้สรุปภาพรวมงานวิจัยด้านการพยากรณ์ราคาวัสดุก่อสร้าง
โดยใช้ข้อมูลเป็นฐาน ชี้ให้เห็นว่าคุณภาพของข้อมูลนำเข้าและกระบวนการเตรียม
ข้อมูลมีผลอย่างยิ่งต่อความถูกต้องของผลลัพธ์การวิเคราะห์ราคา. \textcite{lee2024}
ได้เสนอวิธีการคัดกรอง outlier ด้วยการเปรียบเทียบกับค่าเฉลี่ยและการวิเคราะห์
ค่าความเบี่ยงเบนมาตรฐาน ซึ่งงานวิจัยนี้นำมาประยุกต์ใช้ในขั้นตอน data validation.

จากการวิเคราะห์งานวิจัยที่ผ่านมา ยังขาดงานวิจัยที่บูรณาการกระบวนการทั้งหมด
ตั้งแต่การรวบรวมข้อมูลราคาวัสดุก่อสร้างจากแหล่งภาครัฐและตลาดค้าปลีกสมัยใหม่
การจัดการและปรับมาตรฐานข้อมูลราคา ไปจนถึงการคำนวณต้นทุนวัสดุก่อสร้างต่อหน่วย
ที่สอดคล้องกับหลักเกณฑ์การคำนวณราคากลางของภาครัฐ. งานวิจัยนี้จึงมุ่งเติมเต็ม
ช่องว่างดังกล่าว.
